\begin{textsample}{NEG Dim 4 – 2014\_c5 – Score 1.00 – 2014\_c5/enade\_2014\_ef\_e968405d23.txt}  \label{ex:f4_neg_206}
CAPÍTULO 5 DISTRIBUIÇÃO DOS CONCEITOS 5.1 PANORAMA NACIONAL DA DISTRIBUIÇÃO DOS CONCEITOS O Conceito Enade é um indicador de qualidade que avalia o desempenho dos estudantes a partir dos resultados do ENADE.
Os valores possíveis vão de 1 ( pior situação ) a 5 ( melhor situação ).
As Unidades de Observação com apenas um ou sem nenhum concluinte participante não obtêm o Conceito Enade, ficando Sem Conceito ( SC ).
A Tabela 5.1 apresenta a quantidade e distribuição de cursos de Educação Física ( Licenciatura ) participantes do ENADE/2014, por faixa de conceito e Grande Região.
Mais uma vez, ressaltando que a diferença entre os cursos tabulados neste capítulo e no capítulo 2 são os cursos sem conceito, em princípio, aqueles sem alunos concluintes que participassem da prova ou que tivessem tão somente um aluno concluinte22.
Observando -se os dados da Tabela 5.1, nota -se que, dos 510 cursos participantes, 232 ( 45,5 \% ) classificaram -se com conceito 3, o valor modal.
Este foi também o conceito modal em todas as cinco regiões : Norte ( 42,9 \% ), Nordeste ( 45,5 \% ), Sudeste ( 45,7 \% ), Sul ( 50,5 \% ) e Centro-Oeste ( 35,3 \% ).
O conceito 2 foi o segundo mais frequente em nível nacional ( 26,1 \%, correspondendo a 133 cursos ) e o conceito 4, o terceiro ( 17,6 \%, correspondendo a 90 cursos ).
Houve, ainda, 20 cursos ( 3,9 \% ) que receberam conceito 5 e 24 cursos ( 4,7 \% ) que receberam conceito 1.
Dos 510 cursos de Educação Física ( Licenciatura ), onze cursos ( 2,2 \% ) ficaram sem conceito ( SC ). 22 Estes cursos com somente um concluinte não foram considerados no capítulo 2.
A região Norte participou com 28 cursos ou 5,5 \% do total nacional.
Nesta região, o conceito 3 foi atribuído a 12 cursos dos 28 participantes, o que equivale a 42,9 \% do total regional, além de ser o conceito modal na região.
Aos demais cursos foram atribuídos os conceitos 2 ( 11 cursos ), 4 ( dois cursos ) e 5 ( um curso ).
Dois cursos ficaram sem conceito ( SC ) nessa região.
A região Nordeste participou com 88 cursos ou 17,3 \% do total nacional.
Nesta região, 40 cursos, 45,5 \% em termos regionais, obtiveram conceito 3, o conceito modal para a região.
Os conceitos 2 e 4 foram atribuídos a, respectivamente, 20 e 13 cursos ( 22,7 \% e 14,8 \% ).
Seis cursos obtiveram conceito 1 e sete cursos, conceito 5.
Nessa região, dois cursos ficaram sem conceito ( SC ).
Dos 234 cursos participantes da região Sudeste ( 45,9 \% do total nacional ), o conceito 3 foi o conceito modal, tendo sido obtido por 107 cursos ( 45,7 \% ).
O conceito 1 foi atribuído a 13 cursos ( 5,6 \% ), o conceito 2, a 60 cursos ( 25,6 \% ), o conceito 4, a 42 cursos ( 17,9 \% ) e, por fim, o conceito 5, a sete cursos ( 3,0 \% ).
Cinco cursos ( 2,1 \% ) na região Sudeste ficaram sem conceito ( SC ).
Nos 109 cursos da região Sul, houve predominância do conceito 3, atribuído a 50,5 \% ( conceito modal ), ou seja, a 55 dos 109 cursos participantes na região Sul.
A um único curso ( 0,9 \% ) foi atribuído o conceito 1.
Os demais obtiveram os conceitos 2 ( 25 cursos ), 4 ( 23 cursos ) e 5 ( quatro cursos ), correspondendo, respectivamente, a 22,9 \%, 21,1 \% e 3,7 \% do total da região.
Apenas um curso na região Sul ficou sem conceito ( SC ).
Já dos 51 cursos participantes na região Centro-Oeste, 18 ( 35,3 \% ) receberam conceito 3, o conceito modal, seguidos por outros 17 cursos ( 33,3 \% ), aos quais foi atribuído conceito 2.
Nesta região, somente um curso ficou sem conceito ( SC ), e 4 cursos ( 7,8 \% ) foram avaliados com conceito 1.
Os conceitos 4 e 5 foram atribuídos, respectivamente, a dez cursos ( 19,6 \% ) e a um curso ( 2,0 \% ).
O Gráfico 5.1 apresenta a distribuição do Conceito Enade segundo Grande Região : poligonais mais à esquerda representam Grandes Regiões com pior distribuição, e poligonais mais à direita, Grandes Regiões com melhores conceitos.
A região Norte ( linha azul ) apresenta os cursos em IES com conceitos mais baixos, ao passo que as regiões Sul ( linha roxa ) e Nordeste ( linha verde ) apresentam as distribuições com valores maiores.
A configuração das demais poligonais – regiões Sudeste ( linha vermelha ) e Centro-Oeste ( linha preta ) – apresentam uma situação intermediária. 5.2 CONCEITOS POR CATEGORIA ADMINISTRATIVA E POR GRANDE REGIÃO A Tabela 5.2 apresenta a distribuição dos cursos participantes do ENADE/2014 de Educação Física ( Licenciatura ), por Categoria Administrativa, de acordo com os conceitos por eles alcançados, segundo as Grandes Regiões brasileiras.
Dos 510 cursos participantes, 134 ( 26,3 \% ) eram ministrados em Instituições Públicas e 376 ( 73,7 \% ), em Instituições Privadas.
De acordo com as informações da Tabela 5.2, em termos nacionais, entre Instituições Públicas e Privadas : 20 IES que receberam o conceito 5, dez para cada Categoria Administrativa.
Dos 134 cursos participantes de IES Públicas, o conceito 3 foi o valor modal, atribuído a 57 cursos ( 42,5 \% ).
Entre os demais cursos de Instituições Públicas participantes, oito cursos ( 6,0 \% da categoria ) foram avaliados com conceito 1, 23 cursos ( 17,2 \% ) obtiveram conceito 2, 34 cursos receberam conceito 4 ( 25,4 \% ) e, como já comentado, dez cursos ( 7,5 \% ) foram \textbf{alocados} ao conceito 5.
Nesta categoria, dois cursos ( 1,5 \% ) ficaram sem conceito ( SC ).
Na rede privada, o conceito modal foi 3, atribuído a 175 cursos dos 376 da categoria.
Entre os demais cursos participantes, 16 receberam conceito 1, e 110, conceito 2.
O conceito 4 foi atribuído a 56 cursos e o conceito 5, a dez cursos.
Nesta Categoria Administrativa, nove cursos ficaram sem conceito ( SC ).
Na análise por região, observa -se que, na região Norte, as Instituições Privadas participaram com 17 cursos ( 60,7 \% do total regional ), dos quais apenas um ficou sem conceito ( SC ) e outro curso, com conceito 5.
O conceito modal para as Instituições Privadas na região foi 2, atribuído a nove cursos, correspondendo a 52,9 \%, e seis cursos obtiveram conceito 3.
Nenhum curso foi \textbf{alocado} no conceito 1, nem no conceito 4.
As Instituições Públicas participaram com onze cursos na região Norte ( 39,3 \% do total regional ), dos quais seis obtiveram conceito 3, o valor modal, seguidos por dois cursos ( 18,2 \% ) que receberam conceito 2 e por outros dois cursos que obtiveram conceito 4.
Apenas um curso ficou sem conceito ( SC ), e nenhum curso obteve conceito 1 ou 5 nesta combinação de categoria e região.
Na região Nordeste, a rede privada concentrou 39 dos 88 cursos participantes, o equivalente a 44,3 \% do total da região, a menor das proporções desta categoria entre as regiões.
Dentre os cursos oferecidos por IES Privadas no Nordeste, dois cursos ficaram \textbf{alocados} ao conceito 1, e 14, ao conceito 2.
O conceito 3, valor modal, foi atribuído a 17 cursos.
Quatro cursos receberam conceito 4 e um único curso recebeu conceito 5.
Nesta combinação de categoria e região, apenas um curso ficou sem conceito ( SC ).
As Instituições Públicas da região Nordeste participaram com 49 cursos ( 55,7 \% ), a maior das proporções desta categoria entre as regiões.
Os conceitos 1 e 2 foram atribuídos a quatro e seis cursos, respectivamente.
O valor modal, conceito3, foi obtido por 23 cursos ( 46,9 \% ), seguido por nove cursos que receberam conceito 4 e por outros seis cursos que obtiveram conceito 5.
Somente um curso ficou sem conceito ( SC ) nesta combinação de categoria e região.
Na região Sudeste, a proporção de cursos da rede privada, 85,9 \%, foi mais elevada do que nas demais regiões brasileiras, correspondendo a 201 dos 234 cursos participantes.
Nesta categoria, na região Sudeste, o conceito modal foi 3 ( 91 cursos ).
Nesta combinação de Categoria Administrativa e Grande Região, cinco cursos ficaram sem conceito ( SC ), e outros cinco receberam conceito 5.
Os demais foram avaliados com conceito 1 ( onze cursos ), conceito 2 ( 55 cursos ) e conceito 4 ( 34 cursos ).
Entre os 30 cursos oferecidos por Instituições Públicas na região Sudeste, o conceito modal foi 3 ( 16 cursos ).
Os demais cursos receberam os conceitos 1 ( dois cursos ), 2 ( cinco cursos ), 4 ( oito cursos ) e 5 ( dois cursos ).
Na região Sudeste, nesta categoria, nenhum curso ficou sem conceito ( SC ).
As Instituições Privadas concentraram 84 dos 109 cursos participantes da região Sul, 77,1 \% do total regional.
Desses, 45 ficaram com conceito 3, o conceito modal.
Nesta combinação de Categoria Administrativa e Grande Região, nenhum curso recebeu conceito 1 enquanto 22 cursos receberam conceito 2, 14 cursos obtiveram conceito 4 e dois cursos foram avaliados com conceito 5.
Apenas um curso vinculado a Instituição Privada ficou sem conceito ( SC ) na região Sul.
As Instituições Públicas na região Sul participaram com 25 cursos ( 22,9 \% ).
Os conceitos mais frequentemente atribuídos foram os conceitos 3 ( recebidos por dez cursos, conceito modal ) e 4 ( recebidos por nove cursos ).
Os demais cursos obtiveram os conceitos 1 ( um curso ), 2 ( três cursos ) e 5 ( dois cursos ).
Nenhum curso ficou sem conceito ( SC ) nesta combinação de categoria e região.
Na região Centro-Oeste, 35 dos 51 cursos participantes eram de Instituições Privadas ( 68,6 \% em termos regionais ).
Destes, 16 concentraram -se no conceito 3, conceito modal.
Um curso ficou sem conceito ( SC ).
Quanto aos demais cursos, três receberam conceito 1, dez obtiveram conceito 2, quatro foram avaliados com conceito 4 e apenas a um curso foi atribuído conceito 5.
Dos dezesseis cursos oferecidos por Instituições Públicas, sete foram avaliados com conceito 2, conceito modal, seguidos de seis cursos que obtiveram conceito 4.
Os conceitos 1 e 3 foram atribuídos, respectivamente, a um e a dois cursos.
Nesta região, nenhum curso de IES pública ficou sem conceito ( SC ) ou recebeu conceito 5.
O Gráfico 5.2 apresenta a distribuição do Conceito Enade segundo a Categoria Administrativa da IES.
Os cursos em IES Públicas ( linha azul ) apresentam uma poligonal mais à direita do que os cursos em IES Privadas ( linha verde ) e, consequentemente, uma distribuição de Conceitos Enade com valores maiores. modal, conceito 3 ( recebidos por nove cursos ).
Os demais cursos obtiveram os conceitos 1 ( um curso ), 2 ( três cursos ) e 5 ( dois cursos ).
Nenhum curso ficou sem conceito ( SC ) nesta combinação de categoria e região.
Na região Centro-Oeste, 35 dos 51 cursos participantes eram de Instituições Privadas ( 68,6 \% em termos regionais ).
Destes, 16 concentraram -se no conceito 3, conceito modal.
Um curso ficou sem conceito ( SC ).
Quanto aos demais cursos, três receberam conceito 1, dez obtiveram conceito 2, quatro foram avaliados com conceito 4 e apenas a um curso foi atribuído conceito 5.
Dos dezesseis cursos oferecidos por Instituições Públicas, sete foram avaliados com conceito 2, conceito modal, seguidos de seis cursos que obtiveram conceito 4.
Os conceitos 1 e 3 foram atribuídos, respectivamente, a um e a dois cursos.
Nesta região, nenhum curso de IES pública ficou sem conceito ( SC ) ou recebeu conceito 5.
O Gráfico 5.2 apresenta a distribuição do Conceito Enade segundo a Categoria Administrativa da IES.
Os cursos em IES Públicas ( linha azul ) apresentam uma poligonal mais à direita do que os cursos em IES Privadas ( linha verde ) e, consequentemente, uma distribuição de Conceitos Enade com valores maiores. 5.3 CONCEITOS POR ORGANIZAÇÃO ACADÊMICA E POR GRANDE REGIÃO Na Tabela 5.3, encontra -se a distribuição dos conceitos atribuídos aos cursos participantes do ENADE/2014 na Área de Educação Física ( Licenciatura ), por Organização Acadêmica, segundo as Grandes Regiões brasileiras.
Dos 510 cursos de Educação Física ( Licenciatura ) participantes, 234 eram oferecidos em Universidades, 77 em Centros Universitários e 199 em Faculdades.
Esta distribuição corresponde a, respectivamente, 45,9 \%, 15,1 \% e 39,0 \% dos cursos.
De acordo com os dados apresentados, de todos os 20 cursos avaliados com conceito 5, 14 eram vinculados a Universidades.
Três cursos oferecidos em Universidades ficaram sem conceito ( SC ).
Esse tipo de Organização Acadêmica teve o conceito 3 como modal, com 112 cursos ( 47,9 \% ).
Os demais cursos de Universidades avaliados receberam os conceitos 1 ( sete cursos ), 2 ( 42 cursos ), 4 ( 56 cursos ) e conceito 5 ( 14 cursos, como já mencionado ).
Entre os cursos oferecidos por Centros Universitários, o conceito modal também foi 3, com 38 cursos ( 49,4 \% ).
Neste tipo de Organização Acadêmica, somente um curso ficou sem conceito ( SC ), e os outros cursos receberam os conceitos : 1 ( quatro cursos ), 2 ( 20 cursos ), 4 ( 13 cursos ) e 5 ( um curso ).
Nas Faculdades, sete dos 199 cursos ficaram sem conceito ( SC ), e 82 ( 41,2 \% ) ficaram \textbf{alocados} no conceito 3, conceito modal, seguidos de 71 cursos que obtiveram conceito 2 e de 21 cursos que receberam conceito 4.
Os demais cursos neste tipo de Organização Acadêmica, receberam os conceitos 1 ( 13 cursos ) e 5 ( cinco cursos ).
Considerando -se separadamente as regiões brasileiras, verifica -se que, na região Norte, as Universidades concentraram 12 dos 28 cursos participantes.
Seis cursos ficaram \textbf{alocados} ao conceito 3, conceito modal, seguidos de três cursos que obtiveram conceito 2 e de outros dois cursos que receberam conceito 4.
Apenas um curso ficou sem conceito ( SC ), e nenhum curso recebeu os demais conceitos nesta categoria e região.
Os Centros Universitários da região Norte foram representados por cinco cursos, dos quais três receberam conceito 3 ( conceito modal ) e os outros dois, conceito 2.
As Faculdades participaram com 11 cursos na região Norte que receberam conceitos 2 ( seis cursos, conceito modal ), 3 ( três cursos ) e 5 ( um curso ).
Na região Norte, um curso oferecido em Faculdade ficou sem conceito ( SC ).
Na região Nordeste, as Universidades participaram com 52 dos 88 cursos na Área de Educação Física ( Licenciatura ).
O conceito modal foi 3, atribuído a 26 cursos.
Os demais cursos receberam os conceitos 1 ( quatro cursos ), 2 ( cinco cursos ), 4 ( nove cursos ) e 5 ( seis cursos ).
Dois cursos oferecidos por Universidades no Nordeste ficaram sem conceito ( SC ).
Os Centros Universitários contaram com três cursos participantes na região Nordeste, os quais se distribuíram entre os conceitos 2, 3 e 4 ( um curso cada ).
As Faculdades foram representadas por 33 cursos na região Nordeste, dos quais 14 receberam conceito 2 ( conceito modal ), seguidos de outros 13 que obtiveram conceito 3.
Os conceitos 1, 4 e 5 foram atribuídos, respectivamente, a dois, três e um cursos.
Na região Sudeste, as Universidades concentraram 81 dos 234 cursos oferecidos.
Entre os cursos em Universidades, nessa região, o conceito modal foi 3, atribuído a 41 cursos.
Os demais cursos receberam os conceitos 1 ( um curso ), 2 ( 12 cursos ), 4 ( 22 cursos ) e 5 ( cinco cursos ).
Os Centros Universitários participaram com 51 cursos na região Sudeste, dos quais 24 obtiveram o conceito modal, 3.
Nenhum curso recebeu conceito 5 ou ficou sem conceito ( SC ) nessa categoria e região.
Os demais receberam os conceitos 1 ( quatro cursos ), 2 ( 16 cursos ) e 4 ( sete cursos ).
As Faculdades foram representadas por 102 cursos na região Sudeste.
Nessa categoria e região, oito cursos receberam conceito 1, 32 cursos obtiveram conceito 2, 42 cursos receberam conceito 3 ( conceito modal ), 13 cursos foram avaliados com conceito 4 e a dois cursos foi atribuído conceito 5.
Cinco cursos oferecidos em Faculdades na região Sudeste ficaram sem conceito ( SC ).
Dos 109 cursos da região Sul, 67 eram de Universidades, para os quais o conceito modal foi 3, atribuído a 33 cursos.
Nesse tipo de organização, nenhum dos cursos ficou sem conceito ( SC ).
Os demais receberam os conceitos 1 ( um curso ), 2 ( 14 cursos ), 4 ( 16 cursos ) e 5 ( três cursos ).
Os Centros Universitários da região Sul foram representados por dez cursos.
Exceto por um único curso que ficou sem conceito ( SC ), todos os cursos incluídos nessa categoria e região foram \textbf{alocados} nos conceitos 3 ( seis cursos, conceito modal ) e 4 ( três cursos ).
Já as Faculdades da região Sul foram representadas por 32 cursos, de cujo conjunto nenhum ficou sem conceito ( SC ) ou recebeu conceito 1.
Receberam conceito 3, valor modal, 16 cursos.
Os demais receberam os conceitos 2 ( 11 cursos ), 4 ( quatro cursos ) e 5 ( um curso ).
Na região Centro-Oeste, 22 dos 51 cursos participantes eram de Universidades.
Nesse tipo de organização, nenhum curso ficou sem conceito ( SC ) ou com conceito 5.
O conceito modal foi 2, atribuído a oito cursos.
O segundo conceito mais frequentemente atribuído foi 4, recebido por sete cursos.
Os demais cursos receberam os conceitos 1 ( um curso ) e 3 ( seis cursos ).
Os Centros Universitários da região Centro-Oeste contaram com oito cursos, que obtiveram conceitos 2 ( um curso ), 3 ( quatro cursos, conceito modal ), 4 ( dois cursos ) e 5 ( um curso ).
Nenhum curso ficou sem conceito ( SC ) ou recebeu conceito 1.
Para os 21 cursos oferecidos em Faculdades na região Centro-Oeste, os conceitos modais foram os conceitos 2 e 3, atribuídos, cada qual, a 8 cursos.
Três cursos receberam conceito 1, e um curso obteve conceito 4.
Um curso ficou sem conceito ( SC ) e nenhum recebeu conceito 5.
O Gráfico 5.3 apresenta a distribuição do Conceito Enade segundo a Organização Acadêmica da IES.
Os cursos em Universidades ( linha azul ) apresentam uma poligonal mais à direita do que as demais, denotando uma melhor distribuição de Conceitos Enade.
No outro extremo, temos os cursos oferecidos em Faculdades ( linha vermelha ) com a poligonal mais à esquerda e, consequentemente, uma distribuição de Conceitos Enade com menores valores.
IES de Centros Universitários ( linha verde ) apresentam uma situação intermediária.

% matched lemmas: alocar
\end{textsample}
